\documentclass[11pt,a4paper]{article}
\usepackage[T1]{fontenc}
\usepackage{lmodern,amsmath,amssymb,amsthm,mathtools,booktabs,microtype}
\usepackage[margin=27mm]{geometry}
\usepackage[colorlinks=true,linkcolor=blue,citecolor=blue,urlcolor=blue]{hyperref}
\setlength{\emergencystretch}{2em}
\newtheorem{theorem}{Theorem}[section]
\newtheorem{lemma}[theorem]{Lemma}
\newtheorem{proposition}[theorem]{Proposition}
\newtheorem{corollary}[theorem]{Corollary}
\theoremstyle{remark}\newtheorem{remark}[theorem]{Remark}
\newcommand{\Q}{\mathbb Q}\newcommand{\Qbar}{\overline{\mathbb Q}}
\DeclareMathOperator{\Tr}{Tr}\DeclareMathOperator{\Res}{Res}
\DeclareMathOperator{\divisor}{div}\DeclareMathOperator{\Jac}{Jac}
\title{Native pole equations over abelian parameter spaces}
\author{}\date{16 November 2026}
\begin{document}\maketitle
\begin{abstract}
Fix the input of a native ten-contact or six-contact detector. We prove
that its square-complement equation has no rational section on an
abelian parameter variety when the detector argument is nonconstant.
The original normalized pole equations have no rational witness over
the same function field, even when their square-class representative
and conic parameter vary. At an odd contact divisor, two mandatory
reference tests force the detector input to have negative odd valuation.
Such a divisor must exist: otherwise an etale cover of the parameter
variety would map nontrivially to a curve of genus four or two.
The result concerns uniform function-field constructions. It does not
decide rational lifts after the input, root points, and slope are fixed.
\end{abstract}

\section{Native detectors and the obstruction}

Put
\[
 K=\Q(\sqrt2),\qquad d=3-\sqrt2,\qquad
 E:\ y^2=x^3-d^2x,\qquad A=\Res_{K/\Q}E,
\]
and let $\sigma$ be the nontrivial automorphism of $K$. The source
point and reciprocal abscissa are
\[
 P=\left(\frac{57-28\sqrt2}{8},
          \frac{-660+417\sqrt2}{32}\right),\qquad f=1/x.
\]
On $E$, the function $f$ has divisor $2[O]-2[(0,0)]$.

Let $B/\Q$ be an abelian variety. The point roles of a native source
are rational maps $B\dashrightarrow A$; all source functions and
guards are pulled back along these maps. The input is fixed. For the
general source its early scalar $\rho\in\Q^\times$ is consequently
fixed as well. This includes the scalar
\[
 \rho_*=\frac{102^2\ell^L r}{(1+153\ell^L r)^2}
\]
of \cite{Baseline}. The result uses the unchanged detector factors,
and does not require the arithmetic completeness premises of the source.

\subsection{The ten-contact detector}

For $j=0,1,2$ define
\[
\begin{gathered}
 \Theta=\{0,-1/2,-1,1/2399,-32/407\},\qquad
 \Theta_j=(4j+2+\Theta)\cup(4j+4+\Theta),\\
 \gamma_j=-j-1,\quad
 F_j(T)=\prod_{\theta\in\Theta_j}(T-\theta),\quad
 \delta=t-\gamma_j,\quad w=\rho F_j(t)/\delta^{10}.
\end{gathered}
\]
The ten contacts are distinct and none equals $\gamma_j$.
The native product of \cite[Section~3]{UnitNorm} is
\begin{equation}\label{eq:general}
 \mathcal W=
 \prod_{i=1}^{16}\left(1+\frac{C_i^2}{\delta^2w^{2a_i-1}}\right)
 \prod_{k=1}^{6}\left(1+\frac{G_k^2}{w^{2\kappa_k-1}}\right).
\end{equation}
Here $a_i,\kappa_k$ are positive integers. To specify the required
tests, write $U_0=H$, $V_0=H-P$, and $V_i=U_i-Q_i$ for the other
four root roles. The first ten field tests, in one-based order, are
\[
 (f(U_i))_{i=0}^4,\qquad (\sigma f(V_i))_{i=0}^4.
\]
Their orders are all one. In particular
\begin{equation}\label{eq:mandatory}
             C_1=f(H),\qquad C_6=\sigma f(H-P),
             \qquad a_1=a_6=1.
\end{equation}
The orders of the remaining field tests are one or three. The scalar
factors have exactly the displayed square numerators and odd powers
of $w$; their further formulas are immaterial to the obstruction.

\subsection{The six-contact detector}

For either conditional source of \cite[Section~3]{UnitNorm}, put
\[
 \Theta=\{0,-1/2,-1/4,-2/5,-1,-1/3\},\quad
 F(T)=\prod_{\theta\in\Theta}(T-\theta),\quad
 \delta=t-1,\quad w=\frac{F(t)}{4\delta^8}.
\]
Its product and mandatory tests are
\begin{equation}\label{eq:conditional}
 \mathcal W=\prod_{i=1}^6\left(1+\frac{\delta^4C_i^2}{F(t)}\right),
 \qquad C_1=f(H),\quad C_3=\sigma f(H-P).
\end{equation}
The identity
\[
       \frac{\delta^4C_i^2}{F(t)}
                   =\frac{(C_i/(2\delta^2))^2}{w}
\]
puts every factor in the same form as an order-one factor above.

\subsection{Statement}

In either case write
\begin{equation}\label{eq:input}
        V=\Tr_{K/\Q}(\mathcal W^{-1}),\qquad b=\frac{4w}{V^2}.
\end{equation}
These are the native input formulas. All source denominators and
nonzero guards are assumed nonzero in $\Q(B)$. In particular the
displayed functions are defined, $wV\ne0$, and a constant point role
$H$ is neither $O$ nor $P$ and has the required nonzero abscissas.

\begin{theorem}\label{thm:main}
Suppose the detector argument $t\in\Q(B)$ is nonconstant.
Then
\[
                     V^2-\frac{4w}{17}
\]
is not a square in $\Q(B)$. Equivalently, $17(17-b)$ is not a
square in $\Q(B)$.

More strongly, choose any $d_0,u\in\Q(B)^\times$ with $b=d_0u^2$,
and put $D=17d_0$, $C=17(17-b)$. There are no rational witnesses
over $\Q(B)$ for the normalized pole equations
\begin{equation}\label{eq:poles}
 p^2=D(1-e)(1+3e),\qquad
 q^2=D(s^2-eC)(s^2+3eC).
\end{equation}
This includes all projective conic parameters and the points above
$s=\infty$ on the second-pole curve.
\end{theorem}

The point-role maps may be arbitrary rational maps on $B$. Thus the
theorem permits group translations, isogenies, and combinations of
point roles. It does not fix a pole square class in advance. Its
nonconstant-argument hypothesis is essential to the proof. In the
independent-index families of \cite{Returns}, the input, root points,
and slope are fixed, so every $t$ and $w$ is constant. The theorem
does not settle the rational incidence equations
\[
                      Y_j^2=V_j^2-4w_j/17
\]
on those families, or exclude suitable individual rational tuples.

\section{The two mandatory tests}

\begin{lemma}\label{lem:simple}
The abelian surface $A=\Res_{K/\Q}E$ is simple over $\Q$.
Every nonconstant rational map from an abelian variety over $\Q$
to $A$ extends to a surjective morphism whose geometric fibers have
codimension two.
\end{lemma}
\begin{proof}
At the prime of $K$ above $17$ specified by $\sqrt2=6$, the curve
and its conjugate have good reduction and the following exact counts:
\begin{center}
\begin{tabular}{lrr}\toprule
Curve over $\mathbf F_{17}$ & Points & Frobenius trace\\\midrule
$y^2=x^3-9x$ &20&$-2$\\
$y^2=x^3-13x$&16&$2$\\\bottomrule
\end{tabular}
\end{center}
The discriminants $64(3\mp\sqrt2)^6$ are nonzero at this prime.
A $K$-isogeny would identify the rational $l$-adic Tate modules
for $l\ne17$, hence their good-reduction Frobenius traces.
The unequal traces exclude a $K$-isogeny between $E$ and $E^\sigma$.

If $A$ had an elliptic subvariety over $\Q$, its two projections
after extension to $K$ would both be nonzero: one zero projection
would force the other to be zero by Galois conjugacy. Both projections
would therefore be isogenies, giving a $K$-isogeny between $E$ and
$E^\sigma$. This proves simplicity.

Rational maps from a nonsingular variety to an abelian variety extend
everywhere \cite[Theorem~3.1]{Milne}. A morphism between abelian
varieties is a homomorphism followed by a translation
\cite[Corollary~2.2]{Milne}. Its nonconstant image in $A$ is therefore
all of $A$. In characteristic zero the homomorphism is smooth, and
its fibers are translates of a kernel of codimension two.
\end{proof}

\begin{lemma}[Odd contact divisor]\label{lem:contact}
Let $Z$ be a geometric prime divisor of $B_{\Qbar}$ such that
$v_Z(t-\theta)=m>0$ is odd for one native contact $\theta$.
Then
\begin{equation}\label{eq:valuations}
 v_Z(V)\ge m,\qquad
 v_Z(V^2-4w/17)=m,\qquad
 v_Z(b)=m-2v_Z(V)<0\quad\hbox{is odd}.
\end{equation}
\end{lemma}
\begin{proof}
The contacts are distinct and $\delta$ is a unit at $Z$, so
$v_Z(w)=m$. Each product factor, at either embedding, is of the
form $1+a^2/w^{2r-1}$ up to a square unit at $Z$. For $a\ne0$,
the valuation of its second summand is
\[
                         2v_Z(a)-(2r-1)m,
\]
an odd integer. It cannot be zero or cancel the constant term.
Thus every factor has nonpositive valuation. An identically zero
$a$ gives the factor one and the same conclusion.

Over $\Qbar$ write $A=E_+\times E_-$, $H=(H_+,H_-)$,
and $P=(P_+,P_-)$. At the positive embedding the two mandatory
tests are $f(H_+)$ and $f(H_--P_-)$, and at the negative embedding
they are $f(H_-)$ and $f(H_+-P_+)$.

If $H$ is nonconstant, both positive-embedding tests can have
positive valuation only if $H(Z)=(O,P_-)$. This would place a
divisor in a codimension-two fiber, contrary to
Lemma~\ref{lem:simple}. At least one test therefore has valuation
at most zero. The same argument at $(P_+,O)$ applies to the other
embedding. If $H$ is constant, its generically valid tests are
nonzero algebraic constants and have valuation zero. The argument
is made separately at every geometric prime divisor and does not
require a prime divisor over $\Q$ to remain irreducible.

Each mandatory test has order one. The factor supplied by a test
of nonpositive valuation has valuation at most $-m$; every other
factor has nonpositive valuation. Hence
\[
                 v_Z(\mathcal W_+)\le-m,\qquad
                 v_Z(\mathcal W_-)\le-m.
\]
Their inverses have valuation at least $m$, and so does their sum
$V$, which is generically nonzero. Since $2v_Z(V)>m$, the
summands in $V^2-4w/17$ have unequal valuations. This gives the
second equality in \eqref{eq:valuations}. Formula~\eqref{eq:input}
gives the final assertion.
\end{proof}

\section{Existence of an odd contact divisor}

\begin{lemma}\label{lem:oddexists}
For a nonconstant $t\in\Q(B)$ and either native contact polynomial
$F$, some geometric divisor above a contact has odd multiplicity.
\end{lemma}
\begin{proof}
Suppose otherwise. Every valuation of $F(t)$ is even: this follows
from the supposition at a contact, from
$v_Z(F(t))=(\deg F)v_Z(t)$ at a pole of $t$, and from valuation zero
elsewhere. The degrees $10$ and $6$ are even. Therefore
\[
                            \divisor(F(t))=2\Delta
\]
for an integral Cartier divisor $\Delta$ on the smooth variety
$B_{\Qbar}$. The induced trivialization of
$\mathcal O(2\Delta)$ gives a finite etale double cover with
function algebra obtained by adjoining $\sqrt{F(t)}$. Locally
its equation is $z^2=u$ with $u$ a unit, so it is etale in
characteristic zero. If it splits, take either component.

A connected finite etale cover $B'$ of an abelian variety over an
algebraically closed field is an abelian variety for a suitable
origin, and the cover is an isogeny \cite[Remark~15.3]{Milne}.
The functions $t$ and $\sqrt{F(t)}$ give a nonconstant rational map
from $B'$ to the smooth projective curve
\[
                              C_F:\ y^2=F(t).
\]
Its genus is four or two, since the contacts are distinct.

Embed $C_F$ in its Jacobian using a geometric base point. The
resulting rational map $B'\dashrightarrow\Jac(C_F)$ extends to
a morphism and is a translated homomorphism, by the two results
of \cite{Milne} used above. Its nonconstant image is a closed
abelian translate contained in $C_F$, hence an elliptic curve
equal to $C_F$. This contradicts the genus. An odd contact
divisor must therefore exist.
\end{proof}

Lemmas~\ref{lem:contact} and \ref{lem:oddexists} prove the first
assertion of Theorem~\ref{thm:main}. The square-complement
equivalence follows from
\[
       17(17-b)=\left(\frac{17}{V}\right)^2
                           \left(V^2-\frac{4w}{17}\right).
\]

\section{All normalized pole parameters}

\begin{lemma}\label{lem:oddpole}
Let $F_0$ be a characteristic-zero field with a discrete valuation
$v$ trivial on $\Q^\times$. If $v(b)<0$ is odd, the normalized
pole equations \eqref{eq:poles} have no witness over $F_0$ for
any decomposition $b=d_0u^2$ with $d_0u\ne0$. Both projective
boundaries are included.
\end{lemma}
\begin{proof}
The standard first-conic parametrization is
\[
 e_\tau=\frac{\tau^2-D}{\tau^2+3D},\qquad
 p_\tau=-\frac{4D\tau}{\tau^2+3D}.
\]
Every first-conic point with $e\ne1$ is obtained with
$\tau=p/(e-1)$; the point $(e,p)=(1,0)$ corresponds to
$\tau=\infty$.

Put $k=v(b)<0$. Both $k$ and $v(D)$ are odd. For every finite
$\tau$, the valuation of $\tau^2$ is even or infinite, so it
differs from $v(D)$. The numerator and denominator defining
$e_\tau$ are nonzero and have equal valuation. Thus $v(e_\tau)=0$,
also at $\tau=0$; the parameter infinity gives $e=1$.
Since $v(C)=v(17(17-b))=k$, the quantity $A_\tau=e_\tau C$
has valuation $k$.

For every finite $s$, including zero, the two factors in the
second-pole equation are nonzero and satisfy
\[
 v(s^2-A_\tau)=v(s^2+3A_\tau)=\min\{2v(s),k\}.
\]
Their radicand has odd valuation
\[
 v\bigl(D(s^2-A_\tau)(s^2+3A_\tau)\bigr)
                    =v(D)+2\min\{2v(s),k\}.
\]
It cannot equal $q^2$, including $q=0$. At $s=\infty$, the
chart $u_0=1/s$, $q_0=q/s^2$ requires $q_0^2=D$ at $u_0=0$,
again impossible. The quartic has distinct roots because
$DA_\tau\ne0$ in characteristic zero, so these charts account
for its smooth projective completion.
\end{proof}

Apply Lemma~\ref{lem:oddpole} to the geometric divisor valuation
of Lemma~\ref{lem:contact} over $\Qbar(B)$. This excludes
witnesses over $\Q(B)$ and proves the second assertion of
Theorem~\ref{thm:main}.

\section{Verification and applicability}

The exact checker is
\begin{center}\small
\texttt{negative\_results/computation/abelian\_parameter\_v44/check.py}.
\end{center}
It binds the native source formulas, verifies all contact polynomials
and their distinct roots, independently counts both curves over
$\mathbf F_{17}$, checks the source point and the conditional weight
normalization, and audits the valuation formulas on finite integer
grids. The geometric arguments and universal valuation statements
are proved above; finite enumeration does not replace them.

The obstruction applies to a uniform parametrization by an abelian
variety with nonconstant detector argument. A ramified correspondence
need not admit such a parametrization. Neither suitable rational
specializations nor the fixed-argument incidence problem is excluded.

\begin{thebibliography}{9}
\raggedright
\bibitem{Milne}
J. S. Milne, \emph{Abelian Varieties}, in
G. Cornell and J. H. Silverman (eds.), \emph{Arithmetic Geometry},
Springer, 1986, pp.~103--150.
Author's version dated 2 January 2022:
\url{https://www.jmilne.org/math/xnotes/AVs.pdf}.

\bibitem{UnitNorm}
\emph{A single norm determinant for the 25- and 13-variable constructions},
14 November 2026. Repository source:
\texttt{papers/finite\_algebra/unit\_norm\_20261114\_v28.tex}.

\bibitem{Baseline}
\emph{A scalar normalization for local square-complement realization},
16 November 2026. Repository source:
\texttt{papers/pole\_parameters/}\\
\texttt{square\_complement\_baseline\_20261116\_v43.tex}.

\bibitem{Returns}
\emph{Fixed pole fibers and independent detector indices},
16 November 2026. Repository source:
\texttt{papers/pole\_parameters/}\\
\texttt{common\_parameter\_guard\_reduction\_20261116\_v42.tex}.
\end{thebibliography}
\end{document}
